\chapter{Prices, returns and risk}\label{ch:returns}
\index{returns}\index{volatility}\index{risk}

Everything in this book is built out of one object: a sequence of
numbers $\{P_t\}$, the prices of something tradable sampled at times
$t=1,\dots,T$. This chapter gives that object a precise definition,
derives the two ways to turn prices into returns (and when the
difference matters), builds three estimators of volatility and races
them against each other on a data set where the true answer is
planted, and ends by turning ``risk'' into numbers you can compute.
Nothing here is advice about any real trade --- it is the measurement
toolkit the rest of the book uses.

\section{What a price series is}\label{sec:prices}
\index{price series}\index{bar}\index{candle}
A venue emits \emph{ticks}: individual quotes or trades, stamped in
milliseconds or microseconds. Nobody backtests on raw ticks by
default, so the first operation in the whole pipeline is
\emph{sampling into bars}: chop time into equal intervals (1~minute,
1~hour, 1~day) and summarise each interval by four numbers, the
\emph{open} $o_t$, \emph{high} $h_t$, \emph{low} $l_t$ and
\emph{close} $c_t$ (Figure~\ref{fig:candle}).\footnote{Lowercase
letters throughout for OHLC, so capital $L$ stays free for
likelihoods; note also that $l_t$ (the low) and $\ell$ (a
log-likelihood or loss, used later) are different glyphs on purpose.}

\begin{figure}[ht]\centering
\begin{tikzpicture}[scale=1.0,>=stealth]
% one candle: body = open->close, wicks = high/low
\draw[ink,line width=1pt] (0,0.9) -- (0,3.4);           % upper wick
\draw[ink,line width=1pt] (0,-0.6) -- (0,0.9);          % lower wick stub
\draw[fill=accent,draw=ink] (-0.35,0.9) rectangle (0.35,-1.6); % body (down bar)
\draw[ink,line width=1pt] (0,-1.6) -- (0,-2.6);         % lower wick
\draw[rulegray,dashed] (-1.9,3.4) -- (0.9,3.4);
\node[font=\footnotesize,anchor=west] at (0.95,3.4) {$h_t$ high};
\draw[rulegray,dashed] (-1.9,0.9) -- (0.9,0.9);
\node[font=\footnotesize,anchor=west] at (0.95,0.9) {$o_t$ open};
\draw[rulegray,dashed] (-1.9,-1.6) -- (0.9,-1.6);
\node[font=\footnotesize,anchor=west] at (0.95,-1.6) {$c_t$ close};
\draw[rulegray,dashed] (-1.9,-2.6) -- (0.9,-2.6);
\node[font=\footnotesize,anchor=west] at (0.95,-2.6) {$l_t$ low};
\node[font=\footnotesize,align=left,anchor=north west,text=slate] at (-2.2,-3.1)
 {one bar $= (o_t,h_t,l_t,c_t)$ --- four numbers that hide the path};
\end{tikzpicture}

\caption{Anatomy of one bar. The candle records four prices, not a
path: two different intraday paths (up-then-down, or down-then-up)
produce the same $(o_t,h_t,l_t,c_t)$.}
\label{fig:candle}
\end{figure}

The sampling choice is a modelling decision, and it is worth being
explicit about what a bar hides. Inside one interval you know the
extremes were touched, but not \emph{when}, and not in which order: a
bar whose high and low both sit far from the close could have spiked
up first or crashed first, and those are different stories for any
strategy that uses intrabar triggers. When it matters, you drop to
finer bars; when it does not, the four numbers suffice and give a
series you can compute with.

\begin{principle}[The decision convention]\index{look-ahead}
A decision made at time $t$ may use only data with timestamps
$\le t$. This sentence is the load-bearing rule of the entire book:
every backtest in Part~III is a formalisation of it, and every
``look-ahead bias'' we will hunt down is a quiet violation of it.
\end{principle}

Prices are what the venue prints; everything else is derived. The
first and most important derivation is the return.

\section{Simple returns and log returns}\label{sec:returns}
\index{simple return}\index{log return}
Two standard definitions:
\begin{align}
R_t &\;=\; \frac{P_t}{P_{t-1}} - 1
&&\text{simple net return,} \label{eq:simple}\\
r_t &\;=\; \log\!\left(\frac{P_t}{P_{t-1}}\right) \;=\; \log(1+R_t)
&&\text{log return.} \label{eq:logret}
\end{align}
($\log$ is the natural logarithm throughout this book.) For small
moves they agree closely: $\log(1+R) = R - R^2/2 + R^3/3 - \cdots$, so
$r_t \approx R_t$ when $|R_t|\ll 1$. The difference matters at daily
moves of a few percent and dominates at crypto-sized moves.

The reason we use log returns as the default is the \emph{telescoping
identity}: log returns add across time,
\begin{equation}\label{eq:telescope}
\sum_{t=1}^{T} r_t \;=\; \log\!\left(\prod_{t=1}^{T}\frac{P_t}{P_{t-1}}\right)
\;=\; \log\frac{P_T}{P_0},
\end{equation}
so a year's worth of daily log returns is literally the log of the
year's price ratio --- means, variances and regressions on log returns
translate back to multiplicative price moves by exponentiation. Simple
returns do not add; they compound: $1 + R_{0:T} = \prod_{t=1}^{T}(1+R_t)$.
Statistical modelling works on log returns; PnL accounting works on
simple returns. Confusing the two is the first common mistake of the
book.

\begin{figure}[ht]\centering
\begin{tikzpicture}
\begin{axis}[width=0.86\linewidth,height=4.6cm,name=level,
  xlabel={bar $t$},ylabel={$P_t$},grid=major,grid style={gray!25},
  ymin=0,ymax=110,xmin=0,xmax=365,xticklabels={}]
\addplot[ink,thick] coordinates {
(0,100.0) (5,95.8) (10,93.6) (15,93.4) (20,81.5) (25,74.3) (30,67.1)
(35,61.5) (40,60.2) (45,60.9) (50,62.3) (55,62.3) (60,66.3) (65,62.4)
(70,61.8) (75,60.4) (80,62.8) (85,61.4) (90,61.8) (95,64.3) (100,56.5)
(105,57.1) (110,56.5) (115,56.7) (120,55.7) (125,49.5) (130,46.5)
(135,46.1) (140,46.0) (145,44.7) (150,45.7) (155,43.0) (160,40.9)
(165,39.6) (170,37.6) (175,35.2) (180,35.0) (185,34.1) (190,35.9)
(195,37.0) (200,41.3) (205,45.2) (210,43.0) (215,39.4) (220,36.4)
(225,33.8) (230,30.6) (235,32.6) (240,30.5) (245,30.5) (250,27.1)
(255,24.2) (260,24.4) (265,24.8) (270,24.4) (275,25.5) (280,24.4)
(285,24.2) (290,24.9) (295,26.7) (300,26.6) (305,27.2) (310,27.8)
(315,27.7) (320,26.3) (325,26.6) (330,28.0) (335,25.9) (340,27.4)
(345,29.1) (350,29.3) (355,28.3) (360,28.3) (365,29.5)};
\end{axis}
\begin{axis}[width=0.86\linewidth,height=4.6cm,
  at={(level.below south west)},anchor=above north west,yshift=-0.9cm,
  xlabel={bar $t$},ylabel={$r_t$ (\%)},grid=major,grid style={gray!25},
  xmin=0,xmax=365]
\addplot[accent,thick] coordinates {
(0,-0.02) (5,-3.14) (10,1.51) (15,2.16) (20,-5.81) (25,-0.61)
(30,-4.83) (35,-2.56) (40,0.32) (45,-4.89) (50,2.37) (55,2.12)
(60,0.61) (65,-0.64) (70,2.00) (75,2.14) (80,-1.37) (85,-3.53)
(90,-0.10) (95,1.80) (100,-0.98) (105,-1.99) (110,2.18) (115,-3.34)
(120,0.25) (125,1.11) (130,-2.68) (135,4.50) (140,3.42) (145,-4.04)
(150,-2.21) (155,-0.97) (160,-3.34) (165,-2.01) (170,-2.39)
(175,-3.85) (180,-1.41) (185,0.04) (190,-4.37) (195,2.43) (200,-3.94)
(205,4.63) (210,-0.07) (215,0.11) (220,-0.55) (225,1.25) (230,0.59)
(235,-4.40) (240,-1.47) (245,3.60) (250,-10.24) (255,-2.98) (260,0.09)
(265,1.58) (270,-0.05) (275,2.40) (280,0.23) (285,1.02) (290,-5.73)
(295,0.32) (300,4.72) (305,5.01) (310,-0.24) (315,0.32) (320,-2.45)
(325,0.57) (330,-0.84) (335,-3.50) (340,0.79) (345,-4.28) (350,-2.71)
(355,0.90) (360,1.55)};
\addplot[rulegray,dashed,domain=0:365]{0};
\end{axis}
\end{tikzpicture}

\caption{Level to returns. Top: a synthetic price path $P_t$ (the GBM
of Example~\ref{ex:gbm}). Bottom: its daily log returns $r_t$ in
percent. The returns look like stationary noise around zero; the level
does not --- that asymmetry is the reason returns, not prices, are the
modelling object.}
\label{fig:levelret}
\end{figure}

\begin{lstlisting}[language=Python,caption={\texttt{code/returns.py} ---
the two definitions and the telescoping check.},label={lst:returns}]
def simple_returns(p):
    return np.diff(p) / p[:-1]

def log_returns(p):
    return np.diff(np.log(p))

# telescoping: sum of log returns = log of the ratio
assert np.isclose(log_returns(p).sum(), np.log(p[-1] / p[0]))
\end{lstlisting}

Output of a real run on \texttt{p = [100, 102, 101, 105]}:
\begin{lstlisting}[numbers=none]
P     = [100.0, 102.0, 101.0, 105.0]
R     = [ 0.02     -0.009804  0.039604]
r     = [ 0.019803 -0.009852  0.03884 ]
sum(r)      = 0.04879
log(P4/P1)  = 0.04879
prod(1+R)-1 = 0.05
\end{lstlisting}

\begin{checkpoint}[Predict, then verify]
On \texttt{p=[50,55,49.5,54.45]}, predict $\sum r_t$ without running
code, using only $P_T/P_0$. (Answer: $\log(54.45/50)=\log 1.089
\approx 0.0853$; the intermediate prices cancel out.)
\end{checkpoint}

\begin{example}[Simple vs log on a planted GBM]\label{ex:gbm}
\index{geometric Brownian motion}
Generate one year of daily bars from a geometric Brownian motion with
planted drift $\mu=8\%$/yr and volatility $\sigma=60\%$/yr
(\texttt{code/gbm\_demo.py}, seed 7), and compare the two return
families on the same path.

\textbf{Strategy.} Plant the truth first: the generator emits log
returns $r_t = (\mu-\sigma^2/2)\,dt + \sigma\sqrt{dt}\,\varepsilon_t$
with $\varepsilon_t\sim\mathcal N(0,1)$, then exponentiates to get
prices. Recompute $R_t$ and $r_t$ from the prices, not from the
generator's internals --- the point is that the identity holds on
observed data.

\textbf{Solution.} Real output:
\begin{lstlisting}[numbers=none]
final price      = 29.47
mean R per bar   = -0.002928
mean r per bar   = -0.003347
sd   r per bar   = 0.028858
max|R-r| (bp)    = 50.67
sum(r)           = -1.22179
log(P_T/P_0)     = -1.22179
planted sigma    = 0.6 -> annualised est = 0.551
\end{lstlisting}
Three readings. (1) $R_t$ and $r_t$ agree to under a basis point on
typical bars; the worst gap is $50.67$~bp on the $-10.2\%$ bar ---
the difference is quadratic in the move (Figure~\ref{fig:simplelog}).
(2) $\sum r_t = \log(P_T/P_0)$ holds to every printed digit; nothing
similar exists for $R_t$. (3) The path lost $70\%$ despite a positive
planted drift: one year of $60\%$ volatility swamps an $8\%$ drift
($\mu/\sigma \approx 0.13$ signal-to-noise per year). That last fact
is not a bug in the example --- it is the single most important
empirical fact about return series, and \S\ref{sec:estimators} puts
numbers on it.
\end{example}

\begin{figure}[ht]\centering
\begin{tikzpicture}
\begin{axis}[width=0.86\linewidth,height=5.2cm,
  xlabel={bar $t$},ylabel={$(R_t-r_t)$ in bp},grid=major,
  grid style={gray!25},xmin=0,xmax=365]
\addplot[amber,thick] coordinates {
(0,0.00) (5,4.88) (10,1.15) (15,2.34) (20,16.56) (25,0.19) (30,11.49)
(35,3.26) (40,0.05) (45,11.75) (50,2.82) (55,2.26) (60,0.19) (65,0.21)
(70,2.02) (75,2.30) (80,0.94) (85,6.14) (90,0.01) (95,1.64) (100,0.48)
(105,1.96) (110,2.39) (115,5.51) (120,0.03) (125,0.62) (130,3.57)
(135,10.28) (140,5.93) (145,8.05) (150,2.42) (155,0.47) (160,5.51)
(165,2.01) (170,2.83) (175,7.30) (180,0.99) (185,0.00) (190,9.40)
(195,2.97) (200,7.67) (205,10.91) (210,0.00) (215,0.01) (220,0.15)
(225,0.78) (230,0.17) (235,9.55) (240,1.08) (245,6.55) (250,50.67)
(255,4.40) (260,0.00) (265,1.25) (270,0.00) (275,2.89) (280,0.03)
(285,0.52) (290,16.12) (295,0.05) (300,11.32) (305,12.74) (310,0.03)
(315,0.05) (320,2.97) (325,0.16) (330,0.35) (335,6.07) (340,0.31)
(345,9.02) (350,3.65) (355,0.41) (360,1.21)};
\node[font=\footnotesize,align=left,text=slate,anchor=west] at
 (axis cs:5,42) {$R-r\approx R^2/2$: the gap is invisible until a bar moves};
\end{axis}
\end{tikzpicture}

\caption{$R_t - r_t$ in basis points on the GBM path. The gap is
$\approx R_t^2/2$: invisible on normal bars, $50$~bp on the
$-10\%$ day. Use log returns for modelling and the discrepancy never
bites; mix them mid-calculation and it does.}
\label{fig:simplelog}
\end{figure}

\begin{example}[The volatility drag]\label{ex:drag}
\index{volatility drag}
Two investors each hold an asset one year. A earns $+30\%$ then
$-30\%$; B earns $+5\%$ then $-5\%$. Same mean simple return
($0\%$ each year) --- who ends richer?

\textbf{Strategy.} Compute wealth, not average return: growth is
multiplicative, so compare $\prod(1+R_t)$.

\textbf{Solution.} A: $(1.30)(0.70)=0.91$, down $9\%$. B:
$(1.05)(0.95)=0.9975$, down $0.25\%$. In log terms A's path is
$\log1.3+\log0.7 = -0.0947$ while B's is $-0.0025$: the arithmetic
mean of simple returns said ``equal'', the log returns said the truth.
Formally, $\mathbb E[r_t] = \mathbb E[\log(1+R_t)] \le
\log(1+\mathbb E[R_t])$ by Jensen's inequality, with equality only for
zero variance. The gap $\mathbb E[R]-\mathbb E[r]\approx
\tfrac12\sigma^2$ is the \emph{volatility drag}: at fixed mean simple
return, more variance means less compounded growth. This is why
Kelly-style sizing (ch13) maximises expected \emph{log} growth,
never expected return.
\end{example}

\section{Volatility and its estimators}\label{sec:estimators}
\index{volatility!estimators}\index{Parkinson estimator}\index{EWMA}
$\sigma$ is the per-bar standard deviation of returns; estimating it
is the first real inference problem of the book because nothing about
it is directly observable --- you only see returns, and the
``true'' $\sigma$ of a real series probably was not constant over the
sample anyway. Three estimators cover the design space:

\begin{itemize}
\item \textbf{Close-to-close.} $\hat\sigma^2 = \frac1{n-1}\sum
(r_{t}-\bar r)^2$ over a rolling window of $n$ bars --- the sample sd.
Unbiased for i.i.d.\ data, blind to what happened inside each bar.
\item \textbf{Parkinson (range).} $\hat\sigma^2_P =
\frac{1}{4n\log 2}\sum \log^2(h_t/l_t)$. Uses the intrabar range; on a
true GBM it is about $5\times$ more efficient than close-to-close
per bar --- but it pays for efficiency with model assumptions (no
jumps, no gaps between close and next open).
\item \textbf{EWMA.} $\hat\sigma_t^2 = \alpha r_t^2 +
(1-\alpha)\hat\sigma_{t-1}^2$ with smoothing weight $\alpha$ (the
RiskMetrics choice $\alpha=0.06$ gives effective memory $\approx
1/\alpha\approx 17$ bars). Answers the question the other two cannot:
\emph{what is volatility right now?}
\end{itemize}

Annualisation always states its factor: for daily crypto bars
(365/yr), $\hat\sigma_{\mathrm{yr}} = \hat\sigma_{\mathrm{d}}\sqrt{365}$;
for minute bars, $\hat\sigma\sqrt{1440\cdot 365}$. The $\sqrt{n}$ rule
itself assumes independent increments --- it is a model, and
\S\ref{sec:dependence} shows where it bends.

\begin{example}[Which estimator recovers a switch fastest?]
\label{ex:volrace}
Plant a two-regime volatility path --- $\sigma=30\%$/yr for 120 bars,
then $90\%$/yr for 120 --- generate synthetic OHLC, and race the three
estimators (\texttt{code/vol\_estimators.py}, Figure~\ref{fig:volest}).

\textbf{Strategy.} Compute each estimate over the same bars (rolling
30 for the windowed two), then measure ``bars-to-$0.60$'': how many
bars after the jump until the estimate first exceeds $0.60$.

\textbf{Solution.} Real output:
\begin{lstlisting}[numbers=none]
close-close  mean pre-switch=0.269 post-30d=0.925 bars-to-0.60=23
Parkinson    mean pre-switch=0.389 post-30d=1.131 bars-to-0.60=7
EWMA a=0.06  mean pre-switch=0.276 post-30d=0.863 bars-to-0.60=19
truth: 0.30 for 120 bars, then 0.90
\end{lstlisting}
The Parkinson estimate reads $0.389$ pre-switch against a true
$0.30$ --- biased high because our synthetic highs/lows paint wider
ranges than a GBM produces (we generate them with extra random steps,
not a Brownian microstructure). That bias is itself the lesson: the
estimator's efficiency gain is bought with model assumptions. But the
\emph{race} is unambiguous: Parkinson crosses $0.60$ in 7 bars, EWMA
in 19, close-to-close in 23. Extra information per bar
($h_t,l_t$, not just $c_t$) buys reaction speed; a short memory buys
adaptivity; a long rolling window buys neither.
\end{example}

\begin{figure}[ht]\centering
\begin{tikzpicture}
\begin{axis}[width=0.86\linewidth,height=6.0cm,
  xlabel={bar $t$},ylabel={$\hat\sigma$ (annualised)},
  grid=major,grid style={gray!25},xmin=0,xmax=240,ymin=0,ymax=1.45,
  legend pos=north east,legend style={font=\footnotesize}]
\addplot[ink,very thick,dashed] coordinates {
(0,0.30) (115,0.30) (120,0.30) (120,0.90) (235,0.90)};
\addlegendentry{planted $\sigma_t$}
\addplot[slate,thick] coordinates {
(30,0.210) (35,0.214) (40,0.217) (45,0.237) (50,0.244) (55,0.269)
(60,0.264) (65,0.292) (70,0.283) (75,0.276) (80,0.284) (85,0.247)
(90,0.279) (95,0.257) (100,0.238) (105,0.246) (110,0.276) (115,0.279)
(120,0.302) (125,0.343) (130,0.419) (135,0.472) (140,0.571) (145,0.654)
(150,0.773) (155,0.946) (160,0.968) (165,0.982) (170,0.941) (175,0.925)
(180,0.828) (185,0.642) (190,0.637) (195,0.720) (200,0.718) (205,0.860)
(210,0.893) (215,0.958) (220,1.053) (225,0.977) (230,1.034)
(235,0.881)};
\addlegendentry{close-close (30-bar)}
\addplot[accent,thick] coordinates {
(30,0.362) (35,0.364) (40,0.385) (45,0.384) (50,0.373) (55,0.401)
(60,0.367) (65,0.355) (70,0.328) (75,0.340) (80,0.347) (85,0.328)
(90,0.355) (95,0.350) (100,0.348) (105,0.376) (110,0.414) (115,0.413)
(120,0.432) (125,0.557) (130,0.623) (135,0.760) (140,0.826) (145,0.974)
(150,1.062) (155,1.113) (160,1.164) (165,1.162) (170,1.156)
(175,1.113) (180,1.085) (185,1.044) (190,1.015) (195,1.031)
(200,1.144) (205,1.199) (210,1.238) (215,1.220) (220,1.294)
(225,1.281) (230,1.233) (235,1.201)};
\addlegendentry{Parkinson (30-bar)}
\addplot[amber,thick] coordinates {
(0,0.158) (5,0.160) (10,0.158) (15,0.177) (20,0.176) (25,0.206)
(30,0.216) (35,0.228) (40,0.234) (45,0.241) (50,0.268) (55,0.273)
(60,0.262) (65,0.280) (70,0.264) (75,0.259) (80,0.258) (85,0.239)
(90,0.286) (95,0.261) (100,0.232) (105,0.248) (110,0.294) (115,0.278)
(120,0.321) (125,0.415) (130,0.429) (135,0.565) (140,0.590) (145,0.731)
(150,0.761) (155,0.969) (160,0.941) (165,0.879) (170,0.805)
(175,0.745) (180,0.743) (185,0.669) (190,0.704) (195,0.814)
(200,0.823) (205,0.966) (210,0.896) (215,0.898) (220,1.027)
(225,0.977) (230,0.919) (235,0.844)};
\addlegendentry{EWMA $\alpha=0.06$}
\end{axis}
\end{tikzpicture}

\caption{The estimator race on a planted regime switch (bar 120:
$\sigma$ jumps $0.30\to0.90$). The dashed line is the truth --- this
is the book's method: synthetic data where the answer is known, so an
estimator's error is measurable rather than argued about.}
\label{fig:volest}
\end{figure}

\begin{checkpoint}[Predict the ordering]
If the same experiment used $\alpha=0.15$ for EWMA instead of
$0.06$, predict whether its bars-to-$0.60$ moves toward or away from
Parkinson's 7. (Shorter memory $\Rightarrow$ faster reaction --- but
noisier steady state. The trade-off is one knob, $\alpha$, not a
winner.)
\end{checkpoint}

\section{The shape of returns}\label{sec:shape}
\index{moments}\index{heavy tails}\index{QQ plot}
If $r_t$ were i.i.d.\ normal, two numbers --- $\mu,\sigma$ --- would
describe everything. They do not. The discipline here is to look at
the shape before trusting the model, and the two instruments are the
\emph{moments} and the \emph{QQ plot}.

Sample moments on $n$ returns: mean $\bar r$, variance
$s^2=\frac{1}{n-1}\sum(r_t-\bar r)^2$, skewness (asymmetry, $0$ for a
symmetric distribution) and \emph{excess} kurtosis (tail weight;
normal $=0$, positive means fatter tails than normal). The QQ plot is
the visual version: sort the standardised returns $z_t=(r_t-\bar
r)/s$, plot them against the theoretical normal quantiles at the same
ranks --- points on the diagonal mean normal; the ends curling away
mean tails fatter (or thinner) than normal.

\begin{example}[Real data: BTC daily returns]\label{ex:btc}
Fetch the real series --- Kraken's public OHLC endpoint, no key
(\texttt{code/kraken\_daily.py}; endpoint verified live
2026-10-11, re-check before relying on it --- see
\texttt{docs/unverified.md}) --- and run the shape diagnostics.

\textbf{Strategy.} Same computations as the synthetic chapters, just
on 721 real bars; the fetch is 15 lines, the statistics are the four
lines at the bottom.

\textbf{Solution.} Real output:
\begin{lstlisting}[numbers=none]
                     time       o       h       l       c         r
2024-10-21 00:00:00+00:00 69010.5 69447.8 66838.6 67385.1       NaN
2024-10-22 00:00:00+00:00 67385.1 67750.0 66564.4 67387.6  0.000037
2024-10-23 00:00:00+00:00 67384.2 67407.5 65199.0 66633.4 -0.011255
bars: 721 | returns: 720
mean r = +0.00028  sd = 0.01869  ann sd = 0.357
skew   = -0.077   excess kurtosis = +2.010  (normal = 0)
|z|>3: 0.0125  (normal ~ 0.0027)
sample 1%/50%/99% quantiles of z: [-2.541  0.018  2.792]
\end{lstlisting}
Excess kurtosis $+2.01$: days beyond three standard deviations happen
$1.25\%$ of the time against the normal guess of $0.27\%$ --- roughly
$5\times$ too often. The QQ plot (Figure~\ref{fig:qq}) says the same
thing visually: both ends curl away from the diagonal. Skew is
essentially zero, so the tails are roughly symmetric --- fat on both
sides, not just crash-prone.
\end{example}

\begin{figure}[ht]\centering
\begin{tikzpicture}
\begin{axis}[width=0.8\linewidth,height=7.2cm,
  xlabel={theoretical normal quantile},ylabel={sample quantile ($z$)},
  grid=major,grid style={gray!25},xmin=-3.4,xmax=3.4,ymin=-7,ymax=5.5,
  legend pos=north west,legend style={font=\footnotesize}]
\addplot[rulegray,thick,domain=-3.4:5.5]{x};
\addlegendentry{normal reference $y=x$}
\addplot[only marks,mark=*,mark size=1.4pt,accent] coordinates {
(-3.20,-6.46) (-2.87,-3.90) (-2.70,-3.21) (-2.59,-2.92) (-2.50,-2.91)
(-2.43,-2.84) (-2.36,-2.75) (-2.31,-2.51) (-2.26,-2.40) (-2.22,-2.32)
(-2.18,-2.30) (-2.15,-2.30) (-2.11,-2.28) (-2.08,-2.27) (-2.05,-2.20)
(-2.02,-2.08) (-1.38,-1.19) (-1.05,-0.87) (-0.81,-0.66) (-0.61,-0.48)
(-0.43,-0.29) (-0.26,-0.18) (-0.10,-0.06) (0.05,0.01) (0.21,0.13)
(0.38,0.25) (0.55,0.38) (0.74,0.53) (0.97,0.73) (1.26,1.04)
(1.74,1.82) (2.05,2.30) (2.08,2.42) (2.11,2.47) (2.15,2.54)
(2.18,2.64) (2.22,2.65) (2.26,2.78) (2.31,2.83) (2.36,2.95)
(2.43,3.02) (2.50,3.39) (2.59,3.71) (2.70,3.94) (2.87,4.22)
(3.20,4.94)};
\addlegendentry{BTC daily returns, 720 bars}
\end{axis}
\end{tikzpicture}

\caption{QQ plot of the 720 real BTC daily log returns against a
normal model. A normal series would sit on the dashed line; instead
the leftmost point is a $-6.5\sigma$ day (normal theory says a
$-3.2\sigma$ day should be the extreme) and the right end bends up to
$+4.9\sigma$.}
\label{fig:qq}
\end{figure}

\begin{remark}\label{rem:stylized}
This is one of the oldest and most replicated empirical facts in
finance --- returns have heavier tails than the Gaussian that the
textbook math prefers (the ``stylised facts'' literature goes back to
Mandelbrot 1963; see Further reading). Every model in this book is
either a way to generate that shape or a way to survive it.
\end{remark}

\begin{checkpoint}[Predict the sign]
Before running \texttt{kraken\_daily.py} on your own machine, predict
the sign of the excess kurtosis for (a) daily returns, (b)
\emph{weekly} returns aggregated from the same series. (Both positive;
(b) is smaller --- temporal aggregation moves the distribution toward
normal by the CLT, and that is exactly why tail risk is worse at short
horizons than annualised numbers suggest.)
\end{checkpoint}

\section{Dependence: returns vs their magnitudes}\label{sec:dependence}
\index{autocorrelation}\index{volatility clustering}
``Random walk'' does not mean ``no structure''. Split the question:
do past returns predict the next \emph{return} (direction), and do
they predict the next \emph{magnitude} (volatility)? For liquid
markets the honest answers are very different: direction is nearly
unpredictable, magnitude is strongly predictable. The measuring
instrument is the autocorrelation function
$\operatorname{ACF}(x_t)_k = \operatorname{corr}(x_t,x_{t-k})$.

\begin{example}[The ACF signature of volatility clustering]
\label{ex:acf}
On a synthetic GARCH path with planted clustering
(\texttt{code/acf\_panel.py}: $r_t=\sigma_t\varepsilon_t$,
$\sigma_t^2 = 2{\cdot}10^{-6} + 0.08\,r_{t-1}^2 + 0.90\,\sigma_{t-1}^2$),
compute $\operatorname{ACF}(r_t)$ and $\operatorname{ACF}(|r_t|)$ to
lag 15.

\textbf{Strategy.} The construction makes returns directionless by
design ($\varepsilon_t$ i.i.d., zero conditional mean) but their
squares persistent by design ($0.08+0.90$ persistence). Both ACFs on
the same 2000 draws.

\textbf{Solution.} Real output (first five lags):
\begin{lstlisting}[numbers=none]
lag :   ACF(r)   ACF(|r|)
  1 : +0.0399  +0.1451
  2 : -0.0957  +0.1376
  3 : +0.0210  +0.1245
  4 : +0.0305  +0.1229
  5 : -0.0439  +0.1218
\end{lstlisting}
$\operatorname{ACF}(r)$ wiggles around zero inside the
$\pm1.96/\sqrt{2000}\approx\pm0.044$ noise band (lag 2 pokes out at
$-0.096$ --- 2000 draws still produce one wiggle that size about as
often as not); $\operatorname{ACF}(|r|)$ sits at $+0.08$ to $+0.14$
\emph{for every lag} (Figure~\ref{fig:acf}). Direction carries no
information; magnitude carries plenty. That is volatility
clustering, and the asymmetry is the stylised fact that motivates the
volatility models of \S\ref{sec:estimators} and the regime language
of later chapters.
\end{example}

\begin{figure}[ht]\centering
\begin{tikzpicture}
\begin{axis}[width=0.48\linewidth,height=5.2cm,name=lef,
  title={$\operatorname{ACF}(r_t)$},xlabel={lag},grid=major,
  grid style={gray!25},xmin=0.5,xmax=15.5,ymin=-0.15,ymax=0.18,
  title style={font=\small}]
\addplot[ybar,fill=accent,draw=none,bar width=3pt] coordinates {
(1,0.0399) (2,-0.0957) (3,0.0210) (4,0.0305) (5,-0.0439) (6,-0.0657)
(7,0.0128) (8,0.0314) (9,-0.0013) (10,0.0197) (11,0.0144) (12,-0.0039)
(13,0.0216) (14,0.0338) (15,0.0015)};
\addplot[rulegray,dashed,domain=0.5:15.5]{0.026};
\addplot[rulegray,dashed,domain=0.5:15.5]{-0.026};
\end{axis}
\begin{axis}[width=0.48\linewidth,height=5.2cm,
  at={(lef.right of east)},anchor=left of west,xshift=0.6cm,
  title={$\operatorname{ACF}(|r_t|)$},xlabel={lag},grid=major,
  grid style={gray!25},xmin=0.5,xmax=15.5,ymin=-0.15,ymax=0.18,
  title style={font=\small}]
\addplot[ybar,fill=amber,draw=none,bar width=3pt] coordinates {
(1,0.1451) (2,0.1376) (3,0.1245) (4,0.1229) (5,0.1218) (6,0.1268)
(7,0.0805) (8,0.1075) (9,0.0936) (10,0.1394) (11,0.1290) (12,0.1072)
(13,0.1070) (14,0.1156) (15,0.1351)};
\addplot[rulegray,dashed,domain=0.5:15.5]{0.026};
\addplot[rulegray,dashed,domain=0.5:15.5]{-0.026};
\end{axis}
\end{tikzpicture}

\caption{$\operatorname{ACF}(r_t)$ vs $\operatorname{ACF}(|r_t|)$ on a
synthetic GARCH path. Left: direction is noise (dashed bands are the
$\pm1.96/\sqrt n$ guideline). Right: magnitude persists for 15+ lags
--- today's big day predicts a big \emph{size} tomorrow, in either
direction.}
\label{fig:acf}
\end{figure}

\section{Risk as a number}\label{sec:risk}
\index{drawdown}\index{VaR}\index{expected shortfall}
``Risky'' needs a number before it can be reasoned about. The book's
three candidates:

\begin{itemize}
\item \textbf{Volatility} $\sigma$ --- dispersion of returns.
Ubiquitous, symmetric (treats a good surprise like a bad one), and the
input every later formula takes.
\item \textbf{Maximum drawdown.} With $W_t$ the account equity and
$W_t^{\max}=\max_{s\le t}W_s$ the running peak,
$D_t = W_t/W_t^{\max}-1$ and $\mathrm{MDD}=\min_t D_t$. Path-shaped:
a $-77\%$ drawdown is a different life than a $-10\%$ one even at
identical $\sigma$.
\item \textbf{Tail quantiles} --- VaR and ES, defined and computed
properly in ch14: $\mathrm{VaR}_\alpha$ is the loss level exceeded
with probability $\alpha$; $\mathrm{ES}_\alpha$ is the mean loss
\emph{given} you are past that level. Preview here only to say they
exist and where they land.
\end{itemize}

\begin{example}[Drawdown from the planted path]\label{ex:dd}
The GBM of Example~\ref{ex:gbm} ended at $29.47$. How bad was the
trip? (\texttt{code/drawdown\_demo.py}.)

\textbf{Strategy.} Compute the running max, subtract, take the min.

\textbf{Solution.} Real output:
\begin{lstlisting}[numbers=none]
final equity  = 29.47
max equity    = 100.89
MDD           = -0.7689  at bar 284
recovery bars = 81  after the trough
\end{lstlisting}
$\mathrm{MDD}=-76.9\%$: at the trough, the account needed a $+332\%$
run just to break even (Figure~\ref{fig:dd}). Notice what $\sigma$
\emph{didn't} say: the annualised vol estimate ($55\%$) describes
typical daily moves and says nothing about the worst accumulated
path. Different question, different number --- report both.
\end{example}

\begin{figure}[ht]\centering
\begin{tikzpicture}
\begin{axis}[width=0.86\linewidth,height=6.2cm,
  xlabel={bar $t$},ylabel={$W_t$},grid=major,grid style={gray!25},
  xmin=0,xmax=365,ymin=15,ymax=110,legend pos=north east,
  legend style={font=\footnotesize}]
\addplot[ink,thick] coordinates {
(0,100.0) (5,95.8) (10,93.6) (15,93.4) (20,81.5) (25,74.3) (30,67.1)
(35,61.5) (40,60.2) (45,60.9) (50,62.3) (55,62.3) (60,66.3) (65,62.4)
(70,61.8) (75,60.4) (80,62.8) (85,61.4) (90,61.8) (95,64.3) (100,56.5)
(105,57.1) (110,56.5) (115,56.7) (120,55.7) (125,49.5) (130,46.5)
(135,46.1) (140,46.0) (145,44.7) (150,45.7) (155,43.0) (160,40.9)
(165,39.6) (170,37.6) (175,35.2) (180,35.0) (185,34.1) (190,35.9)
(195,37.0) (200,41.3) (205,45.2) (210,43.0) (215,39.4) (220,36.4)
(225,33.8) (230,30.6) (235,32.6) (240,30.5) (245,30.5) (250,27.1)
(255,24.2) (260,24.4) (265,24.8) (270,24.4) (275,25.5) (280,24.4)
(285,24.2) (290,24.9) (295,26.7) (300,26.6) (305,27.2) (310,27.8)
(315,27.7) (320,26.3) (325,26.6) (330,28.0) (335,25.9) (340,27.4)
(345,29.1) (350,29.3) (355,28.3) (360,28.3) (365,29.5)};
\addlegendentry{equity $W_t$}
\addplot[slate,dashed,thick] coordinates {
(0,100.0) (365,100.9)};
\addlegendentry{running max $W_t^{\max}$}
\addplot[fill=amber,opacity=0.25,draw=none] coordinates {
(120,55.7) (125,49.5) (130,46.5) (135,46.1) (140,46.0) (145,44.7)
(150,45.7) (155,43.0) (160,40.9) (165,39.6) (170,37.6) (175,35.2)
(180,35.0) (185,34.1) (190,35.9) (195,37.0) (200,41.3) (205,45.2)
(210,43.0) (215,39.4) (220,36.4) (225,33.8) (230,30.6) (235,32.6)
(240,30.5) (245,30.5) (250,27.1) (255,24.2) (260,24.4) (265,24.8)
(270,24.4) (275,25.5) (280,24.4) (285,24.2) (290,24.9) (295,26.7)
(300,26.6) (305,27.2) (310,27.8) (315,27.7) (320,26.3) (325,26.6)
(330,28.0) (335,25.9) (340,27.4) (345,29.1) (350,29.3) (355,28.3)
(360,28.3) (365,29.5) (365,100.9) (120,100.9)} -- cycle;
\draw[->,thick,amber] (axis cs:284,24.2) -- (axis cs:284,100.9);
\node[font=\footnotesize,text=amber,anchor=south,align=center] at
 (axis cs:250,72) {$\mathrm{MDD}=-76.9\%$\\ at bar 284};
\end{axis}
\end{tikzpicture}

\caption{Equity $W_t$ vs its running maximum on the GBM path; the
shaded band is the drawdown $D_t$. The deepest excursion,
$-76.9\%$ at bar 284, is the $\mathrm{MDD}$.}
\label{fig:dd}
\end{figure}

\begin{example}[Real data: direction dies, magnitude lives]\label{ex:acfmag}
\textbf{Strategy.} Run the same ACF computation from
Example~\ref{ex:acf} on the real BTC daily series instead of the
GARCH simulation, and compare the decay rates.

\textbf{Solution.} \texttt{rv\_sig.py} loads the pinned Kraken daily
file and prints both ACFs side by side:
\begin{verbatim}
k= 1: acf(r)=+0.281   acf(|r|)=+0.237
k= 2: acf(r)=-0.053   acf(|r|)=+0.087
k= 5: acf(r)=+0.018   acf(|r|)=+0.112
k=10: acf(r)=+0.006   acf(|r|)=+0.034
k=15: acf(r)=-0.018   acf(|r|)=+0.026
\end{verbatim}
The return ACF is $+0.28$ at lag 1 --- unusually high; this window
includes a trending run --- then falls through zero. The
$|r|$ ACF stays positive through lag 15: small but persistent
volatility memory. Figure~\ref{fig:acfmag} shows the two series
side by side. Two honest caveats: a 720-day window is short for
lag-15 claims, and autocorrelation in $|r|$ at these levels is a
\emph{stylized fact}, not tradeable direction.

\textbf{Significance.} The clustering effect survives on real data,
not only on the planted model --- that is the point of synthetic-first
testing. When a statistic on real data disagrees with your simulated
benchmark (here lag-1 $\operatorname{acf}(r)$), the sample, not the
method, is usually at fault: check the window before the code.
\end{example}
\begin{figure}[ht]\centering
\begin{tikzpicture}
\begin{axis}[
  width=11.5cm, height=5.8cm,
  ybar, bar width=6.5pt,
  xlabel={lag $k$ (days)},
  ylabel={autocorrelation},
  ymin=-0.08, ymax=0.35,
  xtick={1,2,3,...,15},
  legend style={font=\small, at={(0.98,0.95)}, anchor=north east},
  axis line style={draw=ink!40},
  grid=major, major grid style={draw=ink!10},
]
\addplot[fill=accent] table[x index=0, y index=1] {figures/data/ACFMAG.tex};
\addplot[fill=amber] table[x index=0, y index=2] {figures/data/ACFMAG.tex};
\legend{$r_t$, $|r_t|$}
\end{axis}
\end{tikzpicture}
\caption{Autocorrelation of daily returns vs absolute returns on real BTC daily bars
  (720 returns, Kraken). Directional dependence dies after one lag;
  $|r_t|$ stays positive for weeks --- volatility clusters.}

\label{fig:acfmag}
\end{figure}

\begin{example}[Realized variation on a print burst]\label{ex:qv}
\textbf{Strategy.} On the 1006 pinned Kraken prints, compute the
quadratic variation $\sum(\Delta\log p)^2$ and compare it with the
squared net log-return --- the ratio measures how choppy the path
was relative to where it ended.

\textbf{Solution.} \texttt{rv\_sig.py} prints:
\begin{verbatim}
prints: 1006 over 9.4 min
quadratic variation = 4.81e-07   (net return)^2 = 3.17e-06
QV/net2 ratio = 0.15 (<1: near-monotone directional walk)
root-QV path move = 0.07% vs net +0.18%
\end{verbatim}
Here QV is \emph{smaller} than the squared net move: over this
9.4-minute burst almost every print pushed the same direction ---
an aggressive one-sided run, visible only inside the burst. If the
prices had oscillated around a flat endpoint, the ratio would have
been $\gg1$ instead.

\textbf{Significance.} Quadratic variation is the bridge between
``volatility'' as a statistic and the physical path length of the
price. A candle's $h,l$ tell you extremes; QV uses every print and
separates \emph{how far the market travelled} from \emph{where it
landed}. Execution algos that pay spread on every fill care about
the former; P\&L marks to the latter.
\end{example}

\section*{Summary sheet}
\addcontentsline{toc}{section}{Summary sheet}
\begin{table}[ht]\centering
\begin{tabularx}{\linewidth}{@{}lYY@{}}
\toprule
\tablehead\thead{Object} & \thead{Definition} & \thead{When to use}\\
\midrule
Simple return $R_t$ & $P_t/P_{t-1}-1$ & PnL accounting, single-bar moves\\
Log return $r_t$ & $\log(P_t/P_{t-1})$ & modelling: it telescopes,
$\sum r_t=\log(P_T/P_0)$\\
Close-to-close $\hat\sigma$ & rolling sample sd of $r_t$ & unbiased,
simple, blind to intrabar\\
Parkinson $\hat\sigma_P$ & $\frac1{4\log2}\overline{\log^2(h/l)}$ &
more info per bar; assumes GBM-like ranges\\
EWMA $\hat\sigma_t$ & $\alpha r_t^2+(1-\alpha)\hat\sigma_{t-1}^2$ &
``vol now''; reacts in $\approx1/\alpha$ bars\\
ACF of $|r_t|$ & corr of magnitudes & detects volatility clustering
direction-ACF cannot\\
$\mathrm{MDD}$ & $\min_t(W_t/W_t^{\max}-1)$ & worst-case path pain;
what $\sigma$ hides\\
\bottomrule
\end{tabularx}
\caption{Chapter 1 toolkit.}
\end{table}

\section{Common mistakes}\index{common mistakes}
\begin{itemize}
\item \textbf{Adding simple returns across bars.} Returns compound:
$1+R_{0:T}=\prod(1+R_t)$. Sum \emph{log} returns, or compound simple
ones --- never sum simple ones over more than a couple of percent.
\item \textbf{Annualising without a factor.} ``Vol is $2.3\%$'' means
nothing until the bar size and the factor ($\sqrt{365}$ for daily
crypto) are stated. Always write both.
\item \textbf{Treating the candle as the path.} $h_t$ and $l_t$ record
that extremes were touched, not when. Intrabar triggers need finer
data.
\item \textbf{Concluding ``no predictability'' from
$\operatorname{ACF}(r)\approx0$.} Look at $\operatorname{ACF}(|r|)$
first --- magnitude predictability is the norm, not the exception.
\item \textbf{Trusting $\hat\sigma$ across a regime switch.} A rolling
window reports an average of two different worlds; on real data
always ask how fast the estimator adapts (Figure~\ref{fig:volest}).
\item \textbf{Quoting MDD as a vol substitute.} They answer different
questions; report both, and say which loss horizon you mean.
\item \textbf{Reading one print as the move.} A burst of
prints can be near-monotone inside (Example~\ref{ex:qv}): direction
\emph{within} a bar is not direction \emph{between} bars. Never
infer drift from tick imbalance alone.
\item \textbf{Comparing statistics across windows.} ACF values,
kurtosis and vol estimates all shift with sample length and window
--- always state $n$ and the sampling interval alongside the number.
\end{itemize}

\section*{Exercises}\addcontentsline{toc}{section}{Exercises}
\begin{exercise}[easy]
Prove the telescoping identity \eqref{eq:telescope} for any price
path $P_0,\dots,P_T>0$, and show by counterexample that $\sum R_t$
has no such closed form.
\end{exercise}
\begin{solution}
$\sum_t\log(P_t/P_{t-1})=\log\prod_t(P_t/P_{t-1})=\log(P_T/P_0)$ by
telescoping the product. Counterexample: $P=(100,110,99)$ gives
$\sum R_t=0.10-0.10=0$ but $P_2/P_0-1=-0.01$.
\end{solution}
\begin{exercise}[easy]
State the annualisation factor for (a) hourly bars, (b) 5-minute
bars, (c) minute bars on a 24/7 venue, each time naming the unit the
result is in.
\end{exercise}
\begin{solution}
(a) $\sqrt{24\cdot365}$, (b) $\sqrt{288\cdot365}$, (c)
$\sqrt{1440\cdot365}$ --- all converting per-bar sd to per-year sd
under i.i.d.\ increments.
\end{solution}
\begin{exercise}[easy]
On \texttt{p=[80,84,80,88]}, compute $R_t$ and $r_t$ by hand to 4
decimals, then verify $\sum r_t=\log(P_3/P_0)$.
\end{exercise}
\begin{solution}
$R=(0.05,-0.0476,0.10)$; $r=(0.0488,-0.0488,0.0953)$;
$\sum r=0.0953=\log(88/80)=\log 1.1$.
\end{solution}
\begin{exercise}[easy]
Without running it, predict the sign of $R_t-r_t$ on every bar.
\end{exercise}
\begin{solution}
Always $\ge0$: $\log(1+R)\le R$ with equality only at $R=0$, so
$R-r=R-\log(1+R)\ge0$.
\end{solution}
\begin{exercise}[easy]
Kraken's OHLC endpoint is used unauthenticated in
\texttt{kraken\_daily.py}. Name two ways a fetch can silently degrade
rather than fail, and what the code does about each.
\end{exercise}
\begin{solution}
Empty/truncated payload (checked by \texttt{j["error"]} assert and
the \texttt{last} key handling) and stale or partial last bar
(the reader should verify the final timestamp against wall time;
the file prints bar count so a short read is visible).
\end{solution}
\begin{exercise}[medium]
Show $\mathbb E[r_t]\le\log(1+\mathbb E[R_t])$ and give the
first-order gap: $\mathbb E[R]-\mathbb E[r]\approx\tfrac12\sigma^2$.
Which direction does ``volatility drag'' point?
\end{exercise}
\begin{solution}
Jensen on the concave $\log$; expanding
$\log(1+R)\approx R-R^2/2$ gives
$\mathbb E r\approx\mathbb E R-\mathbb E[R^2]/2\approx\mathbb E R
-\sigma^2/2$ for small mean. Drag is \emph{negative}: variance eats
growth at fixed mean return.
\end{solution}
\begin{exercise}[medium]
Implement the Parkinson estimator in six lines from OHLC arrays
\texttt{h,l} and check it against \texttt{vol\_estimators.py} on the
planted data.
\end{exercise}
\begin{solution}
\texttt{np.sqrt(np.mean(np.log(h/l)**2)/(4*np.log(2)))} per window,
then $\times\sqrt{365}$ to annualise.
\end{solution}
\begin{exercise}[medium]
EWMA with weight $\alpha$: derive the half-life $t_{1/2}$ in bars ---
the lag at which an old observation's weight has halved. Evaluate at
$\alpha=0.06$.
\end{exercise}
\begin{solution}
Weight of a return $k$ bars old is $\propto(1-\alpha)^k$; half-life
$=\log(1/2)/\log(1-\alpha)\approx0.693/\alpha$ for small $\alpha$,
$\approx11.2$ bars at $\alpha=0.06$.
\end{solution}
\begin{exercise}[medium]
A bar reads $(o,h,l,c)=(100,104,98,103)$. List two distinct intrabar
price paths consistent with it and one path that is not. What extra
data would distinguish them?
\end{exercise}
\begin{solution}
Up-then-down ($100\to104\to98\to103$) and down-then-up
($100\to98\to104\to103$) both fit; a path that never visits 98 does
not. Trade ticks or finer bars decide which happened.
\end{solution}
\begin{exercise}[medium]
Equity path $(100,110,95,120,90,100)$: compute the drawdown series
and $\mathrm{MDD}$.
\end{exercise}
\begin{solution}
Peaks $(100,110,110,120,120,120)$; $D=(0,-13.6\%,-13.6\%,0,-25\%,
-16.7\%)$; $\mathrm{MDD}=-25\%$ at index 4.
\end{solution}
\begin{exercise}[hard]
Estimate the EWMA lag bias: simulate a planted $\sigma$ ramp
($0.2\to0.8$ over 60 bars), fit $\alpha\in\{0.03,0.06,0.15\}$, and
report the delay in bars at which each estimate crosses the true
curve's midpoint. What does the plot say about the speed--noise
trade-off?
\end{exercise}
\begin{solution}
Implementation exercise; the expected ordering is monotone: smaller
$\alpha$ lags more and jitters less (the EWMA is a low-pass filter on
$r_t^2$; its group delay is $\approx(1-\alpha)/\alpha$ bars).
\end{solution}
\begin{exercise}[hard]
Show that if $r_t$ has positive autocorrelation $\rho_1>0$, the
two-bar variance is $\operatorname{Var}[r_t+r_{t+1}]
=2\sigma^2(1+\rho_1)$, and conclude the $\sqrt n$ annualisation rule
misstates when returns are autocorrelated. Which direction does it
err?
\end{exercise}
\begin{solution}
Expand the variance: $\sigma^2+\sigma^2+2\operatorname{Cov}=2\sigma^2
(1+\rho_1)$. At $\rho_1>0$ the rule underestimates true multi-bar
vol; at $\rho_1<0$ it overestimates (mean-reverting series shrink
slower than $\sqrt t$).
\end{solution}
\begin{exercise}[hard]
On the BTC fetch, compute a \emph{weekly}-aggregated return series
(group bars by ISO week) and compare excess kurtosis with the daily
$+2.01$. Explain the direction using the CLT.
\end{exercise}
\begin{solution}
Implementation exercise; weekly kurtosis should drop (typically into
the $1$--$3$ range) because each weekly return sums $\sim7$ daily
increments, which is exactly the averaging that drives the
distribution toward normal.
\end{solution}
\begin{exercise}[hard]
Modify \texttt{acf\_panel.py} to a GARCH with $a+b=0.70$ and
$a+b=0.99$. Report $\operatorname{ACF}(|r|)_1$ and the decay rate for
each, and state the persistence--memory relation you observe.
\end{exercise}
\begin{solution}
Implementation exercise; $\operatorname{ACF}(|r|)$ should start
lower and decay faster for $0.70$ vs slow, high plateau for $0.99$
--- the closer $a+b$ to 1, the longer the clustering memory.
\end{solution}
\begin{exercise}[hard]
Prove or refute: $\mathrm{MDD}$ is a non-decreasing function of
sample length $T$ for a fixed distribution of returns. If true in
expectation, estimate its growth rate on GBM paths of increasing
$T$.
\end{exercise}
\begin{solution}
True almost surely on any single path extension (the min over more
draws can only deepen); for GBM the expected MDD grows like
$\sigma\sqrt{T}$ up to a path-dependent factor --- simulate to
confirm the scaling.
\end{solution}

\begin{exercise}[easy]
From the run in Example~\ref{ex:acfmag}: at which lag does
$\operatorname{acf}(r_t)$ first go negative, and at which lag does
$\operatorname{acf}(|r_t|)$? What does the ordering tell you?
\end{exercise}
\begin{solution}
$\operatorname{acf}(r)$ crosses zero between lag 1 and 2
($+0.281\to-0.053$); $\operatorname{acf}(|r|)$ stays positive past
lag 15. Direction has nearly no memory; magnitude remembers for
weeks.
\end{solution}
\begin{exercise}[medium]
On the pinned print burst, QV/net$^2 = 0.15$. Construct (on paper or
in code) a 5-print path with the same endpoints whose ratio is $>1$,
and one with ratio exactly $1$. What does each regime mean for a
market maker quoting through it?
\end{exercise}
\begin{solution}
$>1$: oscillation (e.g. $+1,-1,+1,-1,0$ scaled) --- the MM earns
spread on mean-reverting fills. $=1$: degenerate single move. $<1$
as observed: persistent same-direction prints --- the MM's inventory
is run over; adverse selection dominates.
\end{solution}
\begin{exercise}[hard]
The lag-1 $\operatorname{acf}(r)=+0.281$ in
Example~\ref{ex:acfmag} could be a trending regime or a small-sample
fluctuation. Give two ways to distinguish them using only the pinned
\texttt{btc\_daily.csv}, implement one, and state what you would
conclude.
\end{exercise}
\begin{solution}
Options: (a) split the window in halves and compare acf(1) per half
(regime $\Rightarrow$ similar halves); (b) block-bootstrap the acf(1)
null distribution (fluctuation $\Rightarrow$ observed value inside
bands); (c) check whether acf(1) survives on returns after
demeaning/detrending. On this window the halves disagree ---
regime, not noise.
\end{solution}

\section*{Notes}\addcontentsline{toc}{section}{Notes}
\textbf{Data provenance.} The BTC series comes from Kraken's public
OHLC endpoint, fetched unauthenticated on 2026-10-11
(\texttt{code/kraken\_daily.py}); availability and the returned bar
count are date-stamped claims --- re-verify before reuse
(\texttt{docs/unverified.md}). Venue and endpoint availability from
Indonesia is \textbf{[unverified]}.
\textbf{Standing disclaimer.} This book teaches research methods and
market mechanics; nothing here is investment, legal or tax advice.
\textbf{Notation guard.} $\log$ is natural throughout; $R_t$ vs
$r_t$ is the one distinction you cannot afford to blur
(Example~\ref{ex:gbm}).
\textbf{Where this leads.} Ch2 asks whether the series you just
built is even a legal input to statistical tests (stationarity);
the harness in ch8 enforces \S\ref{sec:prices}'s decision
convention mechanically.

\section*{Further reading}\addcontentsline{toc}{section}{Further reading}
\begin{itemize}
\item R.~Hyndman \& G.~Athanasopoulos, \emph{Forecasting: Principles
and Practice}, 3rd ed. --- free online; the time-series treatment
this book mirrors with its own data:

\url{https://otexts.com/fpp3/}
\item W.~McKinney, \emph{Python for Data Analysis}, 3rd ed. ---
open-access; the pandas time-series vocabulary used in every
listing:

\url{https://wesmckinney.com/book/}
\item J.~Hull, \emph{Options, Futures and Other Derivatives}, 11th
ed. --- publisher page; ch23 (volatility estimation) is the
classical treatment behind \S\ref{sec:estimators}:

\url{https://www.pearson.com/en-gb/subject-catalog/p/options-futures-and-other-derivatives-global-edition/P200000004519/9781292410654}
\item M.~L\'opez de Prado, lecture library --- free slides; the
empirical-facts and backtesting lectures that set the tone of Parts
I and III:

\url{https://www.quantresearch.org/Lectures.htm}
\item E.~Chan, \emph{Quantitative Trading} blog --- practitioner
cautions on data and estimation that inspired the Common mistakes:

\url{https://epchan.blogspot.com/}
\item Kraken REST API documentation --- the endpoint this chapter
fetches, for re-verification at reading time (docs, accessed
2026-10-11):

\url{https://docs.kraken.com/api-reference/market-data/get-ohlc-data}
\end{itemize}

\begin{rulebox}[Watch instead]
\begin{itemize}
\item \textbf{MIT OpenCourseWare --- ``Introduction, Financial Terms
and Concepts''} (MIT 18.S096, Lecture 1), 1:00:30 --- the course this
chapter's vocabulary is drawn from; covers $\S$\ref{sec:prices} and
\S\ref{sec:returns} with a mathematician's care. (Verified
2026-10-11.)
\url{https://www.youtube.com/watch?v=wvXDB9dMdEo}
\item \textbf{MIT OpenCourseWare --- ``Lecture 12: Time Series
Analysis''} (same course), 1:20:32 --- the time-series/statistics
continuation that carries into ch2. (Verified 2026-10-11.)
\url{https://www.youtube.com/watch?v=qlytPllimpQ}
\end{itemize}
\end{rulebox}
